\section{Justificación}
\label{sec:justificacion}

\subsection{Justificación teórica}

La investigación se fundamenta en la necesidad de modernizar los procesos de gestión escolar mediante la adopción de herramientas tecnológicas contemporáneas. De acuerdo con \textcite{sujovich2024}, la innovación educativa no solo comprende reformas curriculares, sino que <<constituye una transformación práctica de las estructuras institucionales mediante la actualización tecnológica>> (p.~15). La incorporación de software web especializado fortalece la gobernanza institucional y garantiza el principio de equidad en la distribución de oportunidades formativas.

\subsection{Justificación metodológica}

Desde la perspectiva metodológica, el proyecto se respalda en los preceptos de \textcite{hernandez2014}, quienes sostienen que <<la recolección de datos ocurre en los ambientes naturales y cotidianos de los participantes o unidades de análisis>> (p.~397). El levantamiento de información empírica en el propio entorno de San Julián asegura que las optimizaciones técnicas respondan con precisión a las competencias digitales reales de los usuarios locales.

\subsection{Justificación práctica}

En el terreno práctico, la automatización y optimización del registro estudiantil elimina los cuellos de botella administrativos que saturaban la secretaría del establecimiento. Al respecto, \textcite{udb2018} concluyó que <<la automatización del proceso de inscripción reduce considerablemente el tiempo de atención y mejora la eficiencia administrativa>> (p.~27). Al erradicar las visitas presenciales para el reseteo manual de contraseñas, sistematizar la auditoría digital de los requisitos documentales (libretas de calificaciones de 1ro a 3ro de secundaria, cédula de tutores, carnet de identidad y certificado de nacimiento) y automatizar la distribución en cursos y paralelos según el colegio de procedencia, la plataforma web optimizada agiliza el flujo de matrícula y reduce drásticamente las horas de atención extraordinaria del personal del módulo.

\subsection{Justificación técnica}

La solución técnica combina estándares consolidados: semántica HTML5, estilos CSS3 con diseño adaptable \parencite{marcotte2018}, validaciones en cliente con JavaScript \parencite{mozilla2025}, procesamiento en servidor con PHP \parencite{arimetics2024} y un motor relacional en PostgreSQL \parencite{postgresql2025}. Esta selección tecnológica garantiza independencia de licencias propietarias, bajo consumo de ancho de banda y un control transaccional estricto bajo el principio ACID.

\subsection{Justificación social y económica}

En el ámbito social, la propuesta erradica las condiciones degradantes de pernocte en filas y promueve la transparencia institucional, alineándose con el Objetivo de Desarrollo Sostenible (ODS) 4 (Educación Inclusiva y de Calidad) y el ODS 9 (Industria, Innovación e Infraestructura). Económicamente, representa un ahorro familiar colectivo superior a los 20.000 Bs por gestión en pasajes locales (3 Bs por tramo), refrigerios y fotocopias de trámites.
